% GENERATED FILE. Edit canonical lessons, then run npm run generate:books.
% canonical-source-hash: fnv1a64:3df942572bb6dfb0
% canonical-lessons: IT-C429-ostello, IT-C429-agriturismo, IT-C429-locanda, IT-C429-pensione, IT-C429-campeggio

\chapter{Dove dormire}
\label{ch:it-dove-dormire}
\hlchaptermodality{429}

\begin{chapteropening}
\textbf{By the end of this chapter:} \emph{I can choose where to sleep: a hostel, a farm stay, an inn, a guesthouse or a campsite.}

Complete the last lesson of chapter 429: I can choose where to sleep: a hostel, a farm stay, an inn, a guesthouse or a campsite.
\end{chapteropening}

\section[l'ostello]{l'ostello — the hostel}

\label{lesson:IT-C429-ostello}



\begin{quote}
Before the new one: say the Italian for to fasten, then the Italian for the window.
\end{quote}

\subsection*{You'll want to know: l'ostello}
\textbf{l'ostello} — \textquotedblleft{}the hostel\textquotedblright{}.

\begin{quote}
L'ostello costa poco ed è pulito.
\end{quote}

\begin{quote}
\emph{The hostel is cheap and clean.}
\end{quote}

\emph{L'ostello della gioventù} — youth hostel.

\begin{cousinweb}
\emph{ostello} was borrowed from Old French \emph{ostel} 'shelter', which is inherited from Late Latin \emph{hospitāle} 'guesthouse', built on \emph{hospes} 'guest, host' (EtymDB).
\end{cousinweb}

\begin{grammarlens}[title={What you've built}]
One more word for telling, explaining and describing at B1.
\end{grammarlens}

\subsection*{Your turn}
\begin{itemize}
\raggedright
  \item \emph{Say it:} \emph{l'ostello}
  \item \emph{Say it:} \emph{l'ostello}, once more
  \item \emph{Say it:} a sentence of your own with \emph{l'ostello}
  \item \emph{Recall:} say the Italian for to fasten, then the Italian for the window, then say \emph{l'ostello} again
\end{itemize}

\begin{tcolorbox}[breakable,colback=teal!4,colframe=teal!35!black,title={Before you move on}]
What does l'ostello mean? (\textquotedblleft{}The hostel\textquotedblright{}.) Say it once more.
\end{tcolorbox}
\section[l'agriturismo]{l'agriturismo — the farm stay, the farmhouse holiday}

\label{lesson:IT-C429-agriturismo}



\begin{quote}
Before the new one: say the Italian for the window, then the Italian for the hostel.
\end{quote}

\subsection*{You'll want to know: l'agriturismo}
\textbf{l'agriturismo} — \textquotedblleft{}the farm stay, the farmhouse holiday\textquotedblright{}.

\begin{quote}
Abbiamo mangiato benissimo in un agriturismo.
\end{quote}

\begin{quote}
\emph{We ate very well at a farm stay.}
\end{quote}

A working farm that offers rooms and meals, often with its own produce.

\begin{cousinweb}
\emph{agriturismo} is built from \emph{agri-}, from Latin \emph{ager, agri} 'field' (Lewis \& Short), and \emph{turismo}, which was borrowed from English \emph{tourism} (EtymDB).
\end{cousinweb}

\begin{grammarlens}[title={What you've built}]
One more word for telling, explaining and describing at B1.
\end{grammarlens}

\subsection*{Your turn}
\begin{itemize}
\raggedright
  \item \emph{Say it:} \emph{l'agriturismo}
  \item \emph{Say it:} \emph{l'agriturismo}, once more
  \item \emph{Say it:} a sentence of your own with \emph{l'agriturismo}
  \item \emph{Recall:} say the Italian for the window, then the Italian for the hostel, then say \emph{l'agriturismo} again
\end{itemize}

\begin{tcolorbox}[breakable,colback=teal!4,colframe=teal!35!black,title={Before you move on}]
What does l'agriturismo mean? (\textquotedblleft{}The farm stay, the farmhouse holiday\textquotedblright{}.) Say it once more.
\end{tcolorbox}
\section[la locanda]{la locanda — the inn}

\label{lesson:IT-C429-locanda}



\begin{quote}
Before the new one: say the Italian for the hostel, then the Italian for the farm stay.
\end{quote}

\subsection*{You'll want to know: la locanda}
\textbf{la locanda} — \textquotedblleft{}the inn\textquotedblright{}.

\begin{quote}
La locanda è sopra un vecchio ristorante.
\end{quote}

\begin{quote}
\emph{The inn is above an old restaurant.}
\end{quote}

A small, old-fashioned place to eat and sleep.

\begin{cousinweb}
\emph{locanda} comes from Latin \emph{locanda} '(a house) to be let' (EtymDB), a form of \emph{locāre} 'to place, to let', built on \emph{locus} 'place' (Lewis \& Short).
\end{cousinweb}

\begin{grammarlens}[title={What you've built}]
One more word for telling, explaining and describing at B1.
\end{grammarlens}

\subsection*{Your turn}
\begin{itemize}
\raggedright
  \item \emph{Say it:} \emph{la locanda}
  \item \emph{Say it:} \emph{la locanda}, once more
  \item \emph{Say it:} a sentence of your own with \emph{la locanda}
  \item \emph{Recall:} say the Italian for the hostel, then the Italian for the farm stay, then say \emph{la locanda} again
\end{itemize}

\begin{tcolorbox}[breakable,colback=teal!4,colframe=teal!35!black,title={Before you move on}]
What does la locanda mean? (\textquotedblleft{}The inn\textquotedblright{}.) Say it once more.
\end{tcolorbox}
\section[la pensione]{la pensione — the guesthouse; the pension}

\label{lesson:IT-C429-pensione}



\begin{quote}
Before the new one: say the Italian for the farm stay, then the Italian for the inn.
\end{quote}

\subsection*{You'll want to know: la pensione}
\textbf{la pensione} — \textquotedblleft{}the guesthouse; the pension\textquotedblright{}.

\begin{quote}
La pensione è a due passi dal mare.
\end{quote}

\begin{quote}
\emph{The guesthouse is a stone's throw from the sea.}
\end{quote}

\emph{Mezza pensione} — half board. It also means a retirement pension: \emph{andare in pensione}.

\begin{cousinweb}
\emph{pensione} is a learned borrowing from Latin \emph{pēnsiō} 'payment' (Lewis \& Short), from \emph{pendere} 'to weigh out, to pay': first the payment, then the house you paid for.
\end{cousinweb}

\begin{grammarlens}[title={What you've built}]
One more word for telling, explaining and describing at B1.
\end{grammarlens}

\subsection*{Your turn}
\begin{itemize}
\raggedright
  \item \emph{Say it:} \emph{la pensione}
  \item \emph{Say it:} \emph{la pensione}, once more
  \item \emph{Say it:} a sentence of your own with \emph{la pensione}
  \item \emph{Recall:} say the Italian for the farm stay, then the Italian for the inn, then say \emph{la pensione} again
\end{itemize}

\begin{tcolorbox}[breakable,colback=teal!4,colframe=teal!35!black,title={Before you move on}]
What does la pensione mean? (\textquotedblleft{}The guesthouse; the pension\textquotedblright{}.) Say it once more.
\end{tcolorbox}
\section[il campeggio]{il campeggio — the campsite; camping}

\label{lesson:IT-C429-campeggio}



\begin{quote}
Before the new one: say the Italian for the inn, then the Italian for the guesthouse.
\end{quote}

\subsection*{You'll want to know: il campeggio}
\textbf{il campeggio} — \textquotedblleft{}the campsite; camping\textquotedblright{}.

\begin{quote}
Il campeggio ha una piscina.
\end{quote}

\begin{quote}
\emph{The campsite has a pool.}
\end{quote}

\emph{Andare in campeggio} — to go camping.

\begin{cousinweb}
\emph{campeggio} is built on \emph{campo} 'field', which is inherited from Latin \emph{campus} (EtymDB), shaped on the model of English \emph{camping}.
\end{cousinweb}

\begin{grammarlens}[title={What you've built}]
One more word for telling, explaining and describing at B1.
\end{grammarlens}

\subsection*{Your turn}
\begin{itemize}
\raggedright
  \item \emph{Say it:} \emph{il campeggio}
  \item \emph{Say it:} \emph{il campeggio}, once more
  \item \emph{Say it:} a sentence of your own with \emph{il campeggio}
  \item \emph{Recall:} say the Italian for the inn, then the Italian for the guesthouse, then say \emph{il campeggio} again
\end{itemize}

\begin{tcolorbox}[breakable,colback=teal!4,colframe=teal!35!black,title={Before you move on}]
What does il campeggio mean? (\textquotedblleft{}The campsite; camping\textquotedblright{}.) Say it once more.
\end{tcolorbox}
